\section*{Capítulo \ref{ch:g10:why-clouds-form} --- Por qué se forman las nubes}

\begin{solution}{exo:g10:why-clouds-form:1}
La \omterm{def:g10:why-clouds-form:humidity}{humedad relativa} es el contenido de vapor de agua del aire dividido entre
su contenido de \omterm{def:g10:why-clouds-form:humidity}{saturación} a la misma temperatura y presión, en porcentaje.
El \omterm{def:g10:why-clouds-form:dew-point}{punto de rocío} es la temperatura a la que hay que enfriar el aire,
conservando su contenido de agua y su presión, para que se sature.
\end{solution}

\begin{solution}{exo:g10:why-clouds-form:2}
\num{3.8}, \num{7.6}, \num{14.7} y \qty{27.2}{g/kg}. Cada
\qty{10}{\degreeCelsius} lo multiplican aproximadamente por 2
($7.6/3.8 = 2.0$, $14.7/7.6 \approx 1.9$, $27.2/14.7 \approx 1.9$).
\end{solution}

\begin{solution}{exo:g10:why-clouds-form:3}
Contenido de \omterm{def:g10:why-clouds-form:humidity}{saturación} a \qty{24}{\degreeCelsius}: \qty{18.9}{g/kg};
\omterm{def:g10:why-clouds-form:humidity}{humedad relativa} $10.0/18.9 \approx \qty{53}{\percent}$. El contenido de
\omterm{def:g10:why-clouds-form:humidity}{saturación} vale \qty{10.0}{g/kg} a \qty{14}{\degreeCelsius}: es el \omterm{def:g10:why-clouds-form:dew-point}{punto de
rocío}.
\end{solution}

\begin{solution}{exo:g10:why-clouds-form:4}
En verano el aire es húmedo y su \omterm{def:g10:why-clouds-form:dew-point}{punto de rocío} está muy por encima de la
temperatura del vaso: el aire que toca el vaso se enfría por debajo de su
\omterm{def:g10:why-clouds-form:dew-point}{punto de rocío} y el agua se condensa sobre él. El aire interior de invierno,
muy seco, tiene un \omterm{def:g10:why-clouds-form:dew-point}{punto de rocío} inferior a la temperatura del vaso:
enfriado contra él, el aire sigue sin saturarse.
\end{solution}

\begin{solution}{exo:g10:why-clouds-form:5}
Calentamiento desde abajo (cúmulos de la tarde sobre un suelo soleado);
montañas (viento obligado a subir por una cordillera); frentes (aire cálido
que sube sobre aire frío); convergencia (vientos que confluyen cerca del
ecuador o hacia una borrasca).
\end{solution}

\begin{solution}{exo:g10:why-clouds-form:6}
$125 \times (30 - 18) = \qty{1500}{m}$; $125 \times (22 - 20) =
\qty{250}{m}$. El segundo día: las nubes se forman a solo \qty{250}{m} por
encima del suelo.
\end{solution}

\begin{solution}{exo:g10:why-clouds-form:7}
A \qty{1}{km}: $18 - 9.8 = \qty{8.2}{\degreeCelsius}$; a \qty{2}{km}:
$18 - 19.6 = \qty{-1.6}{\degreeCelsius}$. Al bajar, se calienta al mismo
ritmo y vuelve a llegar al suelo a \qty{18}{\degreeCelsius}.
\end{solution}

\begin{solution}{exo:g10:why-clouds-form:8}
Con \qty{4}{\degreeCelsius/km}, tanto una porción no saturada
(\qty{9.8}{\degreeCelsius/km}) como una saturada
(\qty{5}{\degreeCelsius/km}) se enfrían más deprisa que su entorno: \omterm{def:g10:why-clouds-form:stability}{estable}
en ambos casos. Con \qty{12}{\degreeCelsius/km}, incluso una burbuja no
saturada se enfría más despacio que su entorno: sigue más caliente y no deja
de subir, y las burbujas de aire caliente ascienden libremente.
\end{solution}

\begin{solution}{exo:g10:why-clouds-form:9}
$14.7 - 7.6 = \qty{7.1}{g}$ por kilogramo de aire. El \omterm{def:g10:earth-energy-budget:latent-heat}{calor latente} calienta
el aire en el que se condensa el agua, lo que frena su enfriamiento.
\end{solution}

\begin{solution}{exo:g10:why-clouds-form:10}
Los volúmenes crecen como el cubo del diámetro: $(0.5/0.02)^3 = 25^3 = 15\,625$
gotitas para una gota de llovizna; $(2/0.02)^3 = 100^3 = 1\,000\,000$ para
una gota de lluvia.
\end{solution}

\begin{solution}{exo:g10:why-clouds-form:11}
Con el cielo despejado, el suelo radia infrarrojo y se enfría, y con él el
aire que lo toca. Este aire frío es más denso que el de encima; en las
laderas baja pendiente abajo y se acumula en el fondo del valle. Allí sigue
enfriándose durante la noche hasta alcanzar su \omterm{def:g10:why-clouds-form:dew-point}{punto de rocío}, y se forman
gotitas en el propio aire: la niebla. En las cimas de las colinas, el aire
no está tan frío y se mantiene por encima de su \omterm{def:g10:why-clouds-form:dew-point}{punto de rocío}.
\end{solution}

\begin{solution}{exo:g10:why-clouds-form:12}
A la altura $z$, la temperatura es $T - a z$ y el \omterm{def:g10:why-clouds-form:dew-point}{punto de rocío} $T_d - b z$;
son iguales cuando $(a - b) z = T - T_d$, luego $z = (T - T_d)/(a - b)$. Con
$b = 1.7$: $a - b = 8.1$, y $1000/8.1 \approx \qty{123}{m}$ por grado.
\end{solution}

\begin{solution}{exo:g10:why-clouds-form:13}
\omterm{def:g10:why-clouds-form:condensation-level}{Nivel de condensación}: $125 \times (15 - 11) = \qty{500}{m}$; temperatura a
ese nivel: $15 - 9.8 \times 0.5 \approx \qty{10.1}{\degreeCelsius}$. Cresta:
$10.1 - 5
\times 1.5 = \qty{2.6}{\degreeCelsius}$. Bajada de \qty{2}{km}
según el gradiente seco: $2.6 + 9.8 \times 2 \approx \qty{22.2}{\degreeCelsius}$,
unos siete grados más que en la ladera de barlovento.
\end{solution}

\begin{solution}{exo:g10:why-clouds-form:14}
A \qty{-20}{\degreeCelsius}: $126/103 \approx 1.22$, un \qty{22}{\percent}
de sobresaturación respecto al hielo; a \qty{-10}{\degreeCelsius}:
$287/260 \approx
1.10$, un \qty{10}{\percent}. En este intervalo de
temperaturas, cuanto más fría es la nube, mayor es la ventaja del hielo: el
proceso funciona mejor en la parte alta y fría de una nube donde se mezclan
gotitas y cristales.
\end{solution}

\begin{solution}{exo:g10:why-clouds-form:15}
No saturada: la porción se enfría \qty{9.8}{\degreeCelsius} y su entorno
\qty{6.5}{\degreeCelsius}: queda \qty{3.3}{\degreeCelsius} más fría y vuelve
a bajar. Saturada: se enfría \qty{5}{\degreeCelsius}, así que queda
\qty{1.5}{\degreeCelsius} más caliente y sigue subiendo. El aire solo es
\omterm{def:g10:why-clouds-form:stability}{inestable} con una condición: que el aire que sube esté \omterm{def:g10:why-clouds-form:humidity}{saturado}.
\end{solution}

\begin{solution}{pb:g10:why-clouds-form:1}
\textbf{1.} \qty{14.7}{g/kg}.

\textbf{2.} $11.4/14.7 \approx \qty{78}{\percent}$.

\textbf{3.} \qty{16}{\degreeCelsius}, donde el contenido de \omterm{def:g10:why-clouds-form:humidity}{saturación} es de
\qty{11.4}{g/kg}.

\textbf{4.} $11.4/24.1 \approx \qty{47}{\percent}$.

\textbf{5.} No: contiene los mismos \qty{11.4}{g/kg}. Su \omterm{def:g10:why-clouds-form:humidity}{humedad relativa}
bajó porque, más caliente, podía contener más.

\textbf{6.} A su \omterm{def:g10:why-clouds-form:dew-point}{punto de rocío}, \qty{16}{\degreeCelsius}.

\textbf{7.} $11.4 - 8.7 = \qty{2.7}{g}$ por kilogramo de aire.

\textbf{8.} Bajo un cielo nocturno despejado, la hierba y los techos radian
infrarrojo y se enfrían más que el aire; el aire que los toca se enfría por
debajo de su \omterm{def:g10:why-clouds-form:dew-point}{punto de rocío} antes que el resto del aire.

\textbf{9.} $(28.0 - 13.4)/1.5 \approx \qty{9.7}{\degreeCelsius/km}$,
prácticamente el gradiente seco: las burbujas de la tarde han mezclado esta
capa.

\textbf{10.} $(13.4 - 0.4)/2.0 = \qty{6.5}{\degreeCelsius/km}$.

\textbf{11.} \omterm{def:g10:why-clouds-form:stability}{Estable} para una porción no saturada, que se enfría más deprisa
(\qty{9.8}{\degreeCelsius/km}); \omterm{def:g10:why-clouds-form:stability}{inestable} para una saturada, que se enfría
más despacio (\qty{5}{\degreeCelsius/km}).

\textbf{12.} La temperatura aumenta con la altura, de \num{-0.3} a
\qty{1.2}{\degreeCelsius}: una capa muy \omterm{def:g10:why-clouds-form:stability}{estable}, que actúa como una
tapadera. La nube que crece se detiene ahí: sus cimas están cerca de
\qty{4}{km}.

\textbf{13.} $28 - 16 = \qty{12}{\degreeCelsius}$.

\textbf{14.} Temperatura: $28 - 9.8 = \qty{18.2}{\degreeCelsius}$; \omterm{def:g10:why-clouds-form:dew-point}{punto de
rocío}: $16 - 1.8 = \qty{14.2}{\degreeCelsius}$.

\textbf{15.} De 12 a $18.2 - 14.2 = 4.0$: se ha reducido en
\qty{8}{\degreeCelsius}, a un ritmo de \qty{8}{\degreeCelsius} por
kilómetro.

\textbf{16.} Un grado de diferencia se cierra en $1/8$ km, es decir, en
\qty{125}{m}.

\textbf{17.} $125 \times (20 - 16) = \qty{500}{m}$.

\textbf{18.} El \omterm{def:g10:why-clouds-form:dew-point}{punto de rocío} se mantiene en \qty{16}{\degreeCelsius}
mientras el aire cercano al suelo se calienta: la diferencia crece, y con
ella la altura a la que se cierra.

\textbf{19.} $28 - 9.8\,z = 16 - 1.8\,z$.

\textbf{20.} $8\,z = 12$, luego $z = 1.5$: la base de los cúmulos de la
tarde está \textbf{a unos \qty{1.5}{km}} por encima de la llanura (el cálculo
exacto con las propiedades del agua da \qty{1.51}{km}).
\end{solution}
